% GENERATED FILE. Edit canonical lessons, then run npm run generate:books.
% canonical-source-hash: fnv1a64:0086c22b34dfe491
% canonical-lessons: ZH-C227-haibian, ZH-C227-chuan, ZH-C227-yunchuan, ZH-C227-shatan, ZH-C227-yudao

\chapter{A Day at the Sea}
\label{ch:zh-a-day-at-the-sea}
\hlchaptermodality{227}

\begin{chapteropening}
\textbf{By the end of this chapter:} \emph{I can tell about a day at the seaside: the boat, being seasick, the beach and who I ran into.}

Complete the last lesson of chapter 227: I can tell about a day at the seaside: the boat, being seasick, the beach and who I ran into.
\end{chapteropening}

\section[hǎibiān]{hǎibiān — the seaside; the seashore}

\label{lesson:ZH-C227-haibian}



\begin{quote}
Before the new one: say the Mandarin for grassland, then the Mandarin for scenery.
\end{quote}

\subsection*{You'll want to know: hǎibiān}
\emph{hǎibiān} — \textquotedblleft{}the seaside; the seashore\textquotedblright{}.

\begin{quote}
\emph{shàng ge yuè wǒmen qùle hǎibiān} — Last month we went to the seaside.
\end{quote}

Its characters come later; for now it is a word for the ear.

\begin{cousinweb}
\emph{hǎibiān} is built from \emph{hǎi} 'sea' and \zh{边} 'edge', both inherited from Middle Chinese, \emph{xojX} and \emph{pen} (Guangyun, in Baxter's transcription), which explains the first, after an older book, as the pool of heaven that takes in a hundred rivers.
\end{cousinweb}

\begin{grammarlens}[title={What you've built}]
One more word for talking about the past.
\end{grammarlens}

\subsection*{Your turn}
\begin{itemize}
\raggedright
  \item \emph{Say it:} \emph{hǎibiān}
  \item \emph{Say it:} \emph{hǎibiān}, once more
  \item \emph{Say it:} \emph{shàng ge yuè wǒmen qùle hǎibiān} — Last month we went to the seaside.
  \item \emph{Recall:} say the Mandarin for grassland, then the Mandarin for scenery, then say \emph{hǎibiān} again
\end{itemize}

\begin{tcolorbox}[breakable,colback=teal!4,colframe=teal!35!black,title={Before you move on}]
What does hǎibiān mean? (\textquotedblleft{}The seaside; the seashore\textquotedblright{}.) Say it once more.
\end{tcolorbox}
\section[chuán]{chuán — boat; ship}

\label{lesson:ZH-C227-chuan}



\begin{quote}
Before the new one: say the Mandarin for scenery, then the Mandarin for the seaside.
\end{quote}

\subsection*{You'll want to know: chuán}
\emph{chuán} — \textquotedblleft{}boat; ship\textquotedblright{}.

\begin{quote}
\emph{zuótiān wǒmen zuò chuán le} — We went out on a boat yesterday.
\end{quote}

Its characters come later; for now it is a word for the ear.

\begin{cousinweb}
\emph{chuán} is inherited from Middle Chinese \emph{zywen} (Guangyun, in Baxter's transcription), which notes, after an older book, that people west of the passes said this word and people east of them another.
\end{cousinweb}

\begin{grammarlens}[title={What you've built}]
One more word for talking about the past.
\end{grammarlens}

\subsection*{Your turn}
\begin{itemize}
\raggedright
  \item \emph{Say it:} \emph{chuán}
  \item \emph{Say it:} \emph{chuán}, once more
  \item \emph{Say it:} \emph{zuótiān wǒmen zuò chuán le} — We went out on a boat yesterday.
  \item \emph{Recall:} say the Mandarin for scenery, then the Mandarin for the seaside, then say \emph{chuán} again
\end{itemize}

\begin{tcolorbox}[breakable,colback=teal!4,colframe=teal!35!black,title={Before you move on}]
What does chuán mean? (\textquotedblleft{}Boat; ship\textquotedblright{}.) Say it once more.
\end{tcolorbox}
\section[yùnchuán]{yùnchuán — to get seasick}

\label{lesson:ZH-C227-yunchuan}



\begin{quote}
Before the new one: say the Mandarin for the seaside, then the Mandarin for a boat.
\end{quote}

\subsection*{You'll want to know: yùnchuán}
\emph{yùnchuán} — \textquotedblleft{}to get seasick\textquotedblright{}.

\begin{quote}
\emph{wǒ yùnchuán le hěn nánshòu} — I got seasick and felt awful.
\end{quote}

Its characters come later; for now it is a word for the ear.

\begin{cousinweb}
\emph{yùnchuán} is built from \emph{yùn} 'dizzy' and \emph{chuán} 'boat'. The second is inherited from Middle Chinese \emph{zywen} (Guangyun, in Baxter's transcription); the dictionary glosses the first, \emph{hjunH} there, as a halo round the sun or moon, so how it came to mean dizziness is not settled.
\end{cousinweb}

\begin{grammarlens}[title={What you've built}]
One more word for talking about the past.
\end{grammarlens}

\subsection*{Your turn}
\begin{itemize}
\raggedright
  \item \emph{Say it:} \emph{yùnchuán}
  \item \emph{Say it:} \emph{yùnchuán}, once more
  \item \emph{Say it:} \emph{wǒ yùnchuán le hěn nánshòu} — I got seasick and felt awful.
  \item \emph{Recall:} say the Mandarin for the seaside, then the Mandarin for a boat, then say \emph{yùnchuán} again
\end{itemize}

\begin{tcolorbox}[breakable,colback=teal!4,colframe=teal!35!black,title={Before you move on}]
What does yùnchuán mean? (\textquotedblleft{}To get seasick\textquotedblright{}.) Say it once more.
\end{tcolorbox}
\section[shātān]{shātān — sandy beach}

\label{lesson:ZH-C227-shatan}



\begin{quote}
Before the new one: say the Mandarin for a boat, then the Mandarin for to get seasick.
\end{quote}

\subsection*{You'll want to know: shātān}
\emph{shātān} — \textquotedblleft{}sandy beach\textquotedblright{}.

\begin{quote}
\emph{nàr de shātān hěn gānjìng} — The beach there was really clean.
\end{quote}

Its characters come later; for now it is a word for the ear.

\begin{cousinweb}
\emph{shātān} is built from \emph{shā} 'sand' and \emph{tān} 'a shore', both inherited from Middle Chinese, \emph{srae} and \emph{than} (Guangyun, in Baxter's transcription), which explains the first, after an older book, as stones scattered by water, and glosses the second as shallows in a river, so how it came to mean a beach is not settled.
\end{cousinweb}

\begin{grammarlens}[title={What you've built}]
One more word for talking about the past.
\end{grammarlens}

\subsection*{Your turn}
\begin{itemize}
\raggedright
  \item \emph{Say it:} \emph{shātān}
  \item \emph{Say it:} \emph{shātān}, once more
  \item \emph{Say it:} \emph{nàr de shātān hěn gānjìng} — The beach there was really clean.
  \item \emph{Recall:} say the Mandarin for a boat, then the Mandarin for to get seasick, then say \emph{shātān} again
\end{itemize}

\begin{tcolorbox}[breakable,colback=teal!4,colframe=teal!35!black,title={Before you move on}]
What does shātān mean? (\textquotedblleft{}Sandy beach\textquotedblright{}.) Say it once more.
\end{tcolorbox}
\section[yùdào]{yùdào — to run into; to come across}

\label{lesson:ZH-C227-yudao}



\begin{quote}
Before the new one: say the Mandarin for to get seasick, then the Mandarin for a sandy beach.
\end{quote}

\subsection*{You'll want to know: yùdào}
\emph{yùdào} — \textquotedblleft{}to run into; to come across\textquotedblright{}.

\begin{quote}
\emph{zuótiān wǒ yùdàole yí ge lǎo tóngxué} — Yesterday I ran into an old classmate.
\end{quote}

Its characters come later; for now it is a word for the ear.

\begin{cousinweb}
\emph{yùdào} is built from \emph{yù} 'to meet by chance' and \zh{到} 'to arrive', both inherited from Middle Chinese, \emph{ngjuH} and \emph{tawH} (Guangyun, in Baxter's transcription), which glosses the first 'to meet without arranging it'.
\end{cousinweb}

\begin{grammarlens}[title={What you've built}]
One more word for talking about the past.
\end{grammarlens}

\subsection*{Your turn}
\begin{itemize}
\raggedright
  \item \emph{Say it:} \emph{yùdào}
  \item \emph{Say it:} \emph{yùdào}, once more
  \item \emph{Say it:} \emph{zuótiān wǒ yùdàole yí ge lǎo tóngxué} — Yesterday I ran into an old classmate.
  \item \emph{Recall:} say the Mandarin for to get seasick, then the Mandarin for a sandy beach, then say \emph{yùdào} again
\end{itemize}

\begin{tcolorbox}[breakable,colback=teal!4,colframe=teal!35!black,title={Before you move on}]
What does yùdào mean? (\textquotedblleft{}To run into; to come across\textquotedblright{}.) Say it once more.
\end{tcolorbox}
